% =============================================================================
% Semana 05 — Histogramas, Distribuições e Dados Ausentes
% Aula de 3 horas (19h–22h) | 11/03/2026
% =============================================================================

\chapter{Semana 05 (11/03/2026) — Histogramas, Distribuições e Dados Ausentes}

% ---------------------------------------------------------------------------
\section{Objetivo da semana}
% ---------------------------------------------------------------------------

Pessoal, nas últimas duas semanas calculamos \emph{onde} os dados estão
(posição) e \emph{quão espalhados} eles estão (dispersão). Hoje vamos olhar
para o \textbf{formato} dos dados: o histograma. Com ele respondemos perguntas
como — ``o comportamento é simétrico ou enviesado?'', ``há uma cauda
longa de valores extremos?'', ``a distribuição tem um pico ou dois?''

Depois enfrentamos o problema mais frequente na vida real: \textbf{dados
faltando}. Missing values não são só uma inconveniência técnica — a ausência
de um dado carrega informação e a decisão de como tratá-la afeta completamente
as análises subsequentes.

\begin{BoardBox}
\textbf{Roteiro da aula (19h–22h)}
\begin{enumerate}
  \item 19h–19h25 — \textbf{REVISÃO NO QUADRO:} Média, mediana, variância, outlier, covariância, correlação
  \item 19h25–19h45 — \textbf{TEORIA:} O que é um histograma e como ler bins
  \item 19h45–20h10 — \textbf{TEORIA:} Assimetria, caudas e distribuição normal
  \item 20h10–20h30 — \textbf{PRÁTICA EM DUPLA:} Tabela 1 — construir e interpretar histograma
  \item 20h30–21h — \textbf{TEORIA:} Dados ausentes — tipos, diagnóstico e estratégias
  \item 21h–21h30 — \textbf{PRÁTICA EM DUPLA:} Tabela 2 — diagnóstico e tratamento de missing
  \item 21h30–21h45 — \textbf{TEORIA:} Preparação final para Tableau
  \item 21h45–22h — \textbf{ATIVIDADE INDIVIDUAL FINAL}
\end{enumerate}
\end{BoardBox}

\begin{NoteBox}
\textbf{Leitura base:} Cap.\ 02 (seções de histograma e distribuição) + Cap.\ 01
(seção de dados ausentes). Livro: \textit{Estatística Prática para Cientistas de Dados} (Bruce).
\end{NoteBox}

% =============================================================================
\section{Revisão no Quadro — Conteúdos das Semanas 01–04}
% =============================================================================

\begin{NoteBox}
Antes do conteúdo novo, vamos revisar os fundamentos das semanas
anteriores com questões resolvidas no quadro. Foco: média, mediana,
moda, variância, desvio padrão, outlier, covariância e correlação.
\end{NoteBox}

\begin{SolvedBox}
\textbf{Q1 — Média, Mediana e Moda}

Uma loja de departamentos registrou vendas diárias (R\$ mil) em 7 dias:\\
$\{12,\; 15,\; 14,\; 13,\; 60,\; 15,\; 16\}$

\begin{enumerate}
  \item Calcule média, mediana e moda.
  \item Qual medida é mais representativa? Por quê?
\end{enumerate}

\textbf{Resolução:}

\textbf{Média:} $\bar{x} = \frac{12+15+14+13+60+15+16}{7} = \frac{145}{7} \approx 20{,}7$

\textbf{Ordenados:} $12, 13, 14, \mathbf{15}, 15, 16, 60$ \quad
$\Rightarrow$ \textbf{Mediana} $= 15$ \quad (4\textsuperscript{a} posição, $n=7$)

\textbf{Moda:} 15 (aparece 2$\times$)

A \textbf{mediana} (R\$ 15 mil) é mais representativa. O valor atípico de
R\$ 60 mil (provavelmente uma promoção) puxa a média para cima — a média
de R\$ 20,7 mil não reflete nenhum dia ``normal''.
\end{SolvedBox}

\begin{SolvedBox}
\textbf{Q2 — Variância, Desvio Padrão e Outlier}

Duas filiais de uma franquia registraram vendas (R\$ mil) em 5 dias:

\begin{center}
\begin{tabular}{|l|r|r|r|r|r|c|}
\hline
& Dia 1 & Dia 2 & Dia 3 & Dia 4 & Dia 5 & $\bar{x}$ \\
\hline
Filial A & 10 & 11 & 10 & 12 & 12 & 11 \\
Filial B & 5 & 20 & 8 & 15 & 7 & 11 \\
\hline
\end{tabular}
\end{center}

(a) Calcule o desvio padrão de cada filial.\\
(b) Qual filial é mais previsível?\\
(c) Para a Filial B, com $Q_1 = 6{,}0$ e $Q_3 = 17{,}5$: o valor 20 é outlier?

\textbf{Resolução:}

\textbf{Filial A} ($\bar{x}=11$): desvios$^2 = 1+0+1+1+1 = 4$ \quad
$s_A^2 = 4/4 = 1{,}0$ \quad $s_A = 1{,}0$

\textbf{Filial B} ($\bar{x}=11$): desvios$^2 = 36+81+9+16+16 = 158$ \quad
$s_B^2 = 158/4 = 39{,}5$ \quad $s_B \approx 6{,}3$

\textbf{(b)} Filial A ($s=1{,}0$) é muito mais previsível. Mesma média,
comportamentos completamente diferentes.

\textbf{(c)} $IQR = 17{,}5 - 6{,}0 = 11{,}5$ \quad
$LS = Q_3 + 1{,}5 \times IQR = 17{,}5 + 17{,}25 = 34{,}75$\\
$20 < 34{,}75$ $\Rightarrow$ 20 \textbf{não é outlier} pela regra de Tukey.
\end{SolvedBox}

\begin{SolvedBox}
\textbf{Q3 — Covariância e Correlação de Pearson}

Uma equipe mediu horas de treinamento ($X$) e unidades produzidas ($Y$)
em 4 meses:

\begin{center}
\begin{tabular}{|c|r|r|}
\hline
\textbf{Mês} & \textbf{Treino ($X$)} & \textbf{Produção ($Y$)} \\
\hline
Jan & 3 & 20 \\
Fev & 5 & 30 \\
Mar & 7 & 50 \\
Abr & 9 & 40 \\
\hline
\end{tabular}
\end{center}

(a) Calcule $\bar{x}$, $\bar{y}$ e a covariância.\\
(b) Calcule $r$ (Pearson) e interprete.\\
(c) Podemos afirmar que mais treinamento \emph{causa} mais produção?

\textbf{Resolução:}

$\bar{x} = 6$, \quad $\bar{y} = 35$

\begin{center}
\begin{tabular}{|r|r|r|r|c|}
\hline
$(x_i{-}6)$ & $(y_i{-}35)$ & $(x{-}\bar{x})^2$ & $(y{-}\bar{y})^2$ & Produto \\
\hline
$-3$ & $-15$ & 9 & 225 & 45 \\
$-1$ & $-5$  & 1 & 25  & 5  \\
$+1$ & $+15$ & 1 & 225 & 15 \\
$+3$ & $+5$  & 9 & 25  & 15 \\
\hline
\multicolumn{2}{|r|}{\textbf{Soma}} & 20 & 500 & \textbf{80} \\
\hline
\end{tabular}
\end{center}

\textbf{(a)} $\text{Cov}(X,Y) = \frac{80}{3} \approx 26{,}7$ (positiva: sobem juntas)

\textbf{(b)} $s_X = \sqrt{20/3} \approx 2{,}58$; \quad
$s_Y = \sqrt{500/3} \approx 12{,}91$

\[
r = \frac{26{,}7}{2{,}58 \times 12{,}91} = \frac{26{,}7}{33{,}3} \approx 0{,}80
\]

Correlação positiva \textbf{forte} — meses com mais horas de treinamento
estão associados a maior produção.

\textbf{(c)} Não. Correlação $\neq$ causalidade. A produção pode ter subido
por outros fatores (sazonalidade, motivação, novos equipamentos). Para
afirmar causalidade, seria preciso um experimento controlado.
\end{SolvedBox}

% =============================================================================
\section{Parte I — Histogramas e Distribuições (19h–20h)}
% =============================================================================

% ---------------------------------------------------------------------------
\subsection{19h–19h20: O que é um histograma e como ler bins}
% ---------------------------------------------------------------------------

Pessoal, o histograma é o gráfico mais honesto que existe. Ele mostra
\textbf{exatamente como os dados se distribuem} — sem suavizar, sem resumir
em um único número.

\begin{BoardBox}
\textbf{Definição e anatomia do histograma (para o quadro)}

\textbf{O que é:} o histograma agrupa os dados em intervalos (\textbf{bins})
e mostra quantas observações caem em cada intervalo. O eixo X é o
valor da variável; o eixo Y é a contagem (ou frequência relativa).

\textbf{Por que é diferente de um gráfico de barras?}
\begin{itemize}
  \item Barras: variável categórica (A, B, C). Barras não se tocam.
  \item Histograma: variável numérica contínua. Barras se tocam — porque
  o intervalo de um bin começa onde o anterior termina.
\end{itemize}

\textbf{Escolha do número de bins:}
\begin{itemize}
  \item Poucos bins: esconde padrões, tudo parece uniforme.
  \item Muitos bins: ruído visual, impossível ver padrão.
  \item Regra de Sturges: $k = 1 + 3{,}3 \cdot \log_{10}(n)$ — ponto de partida.
  \item Na prática: experimente 10–20 bins e ajuste pelo contexto.
\end{itemize}

\textbf{Exemplo:} 50 salários entre R\$1.000 e R\$12.000, com bins de R\$1.000:\\
Bin [1000–2000): 18 funcionários\\
Bin [2000–3000): 14 funcionários\\
Bin [3000–4000): 9 funcionários\\
$\ldots$ (poucos acima de R\$8.000)\\
Leitura: maioria concentrada na faixa de menor renda.
\end{BoardBox}

% ---------------------------------------------------------------------------
\subsection{19h20–20h: Assimetria, caudas e distribuição normal}
% ---------------------------------------------------------------------------

\begin{SolvedBox}
\textbf{Formas de distribuição — o vocabulário essencial:}

\textbf{1. Distribuição simétrica / Normal (Gaussiana):}
\begin{itemize}
  \item Forma de sino. Média $\approx$ Mediana $\approx$ Moda.
  \item A maioria dos dados concentra-se no centro.
  \item \textbf{Regra empírica (68-95-99,7):} em uma distribuição normal,
  aproximadamente 68\% dos dados caem dentro de $\pm 1\sigma$, 95\% dentro
  de $\pm 2\sigma$, e 99,7\% dentro de $\pm 3\sigma$.
\end{itemize}

\textbf{2. Assimetria positiva (cauda à direita):}
\begin{itemize}
  \item A maioria dos dados está à esquerda; há uma cauda de valores altos.
  \item Média $>$ Mediana. Ex.: renda, preço de imóveis, tempo de resolução de bugs.
  \item A média é ``puxada'' pelos extremos: usar mediana é mais representativo.
\end{itemize}

\textbf{3. Assimetria negativa (cauda à esquerda):}
\begin{itemize}
  \item Maioria à direita; há uma cauda de valores baixos.
  \item Média $<$ Mediana. Ex.: notas de prova difícil, tempo até falha de equipamento.
\end{itemize}

\textbf{4. Distribuição bimodal:}
\begin{itemize}
  \item Dois picos — indica a existência de dois subgrupos distintos.
  \item Ex.: pesquisar todos os salários de uma empresa que mistura estagiários
  e sêniors. A análise correta exige separar os grupos.
\end{itemize}
\end{SolvedBox}

\begin{FormulaBox}{Coeficiente de assimetria (Skewness)}
\[
g_1 = \frac{\frac{1}{n}\sum_{i=1}^{n}(x_i - \bar{x})^3}{s^3}
\]
\begin{itemize}
  \item $g_1 > 0$: assimetria positiva (cauda direita)
  \item $g_1 = 0$: distribuição simétrica
  \item $g_1 < 0$: assimetria negativa (cauda esquerda)
\end{itemize}
Na prática, o valor importa menos do que a direção. Valores $|g_1| > 1$
já indicam assimetria relevante.
\end{FormulaBox}

\begin{BoardBox}
\textbf{Regra de ouro da assimetria + escolha de medida central (para o quadro)}
\begin{tabular}{|l|l|l|}
\hline
\textbf{Distribuição} & \textbf{Use como centro} & \textbf{Use como dispersão} \\
\hline
Simétrica, sem outliers & Média & Desvio Padrão \\
Assimétrica ou com outliers & Mediana & IQR \\
Bimodal & Descreva os dois grupos & DP de cada grupo \\
\hline
\end{tabular}
\end{BoardBox}

Perguntem aos alunos: \textit{``O histograma de renda de um país geralmente
tem assimetria positiva intensa. O que isso significa sobre a representatividade
de usar a `renda média' como indicador oficial? Que número seria mais honesto?''}

\begin{SolvedBox}
\textbf{Questão para o Quadro — Identificando a Forma da Distribuição}

O histograma de valores de empréstimos de um banco digital mostra:

\begin{center}
\begin{tabular}{|l|r|}
\hline
\textbf{Faixa (R\$ mil)} & \textbf{Contratos} \\
\hline
1–5     & 450 \\
5–10    & 280 \\
10–20   & 150 \\
20–50   & 80  \\
50–200  & 40  \\
\hline
\end{tabular}
\end{center}

(a) Que forma de distribuição é essa?\\
(b) A média será maior ou menor que a mediana?\\
(c) Que medida de centro você reportaria ao regulador do Banco Central?

\textbf{Resolução:}\\
(a) \textbf{Assimetria positiva} — concentração na faixa baixa, cauda longa à direita.\\
(b) \textbf{Média $>$ Mediana} — a cauda de empréstimos altos puxa a média para cima.\\
(c) \textbf{Mediana} — representa melhor o contrato ``típico''. A média seria
inflada pelos poucos empréstimos de alto valor.
\end{SolvedBox}

% =============================================================================
\section{Parte II — Prática em Dupla 1 (20h–20h30)}
% =============================================================================

\begin{BoardBox}
\textbf{Tabela 1 — Tempo de Resolução de Chamados de TI (em minutos)}

\medskip
\begin{tabular}{|c|l|r|l|}
\hline
\textbf{ID} & \textbf{Chamado} & \textbf{Tempo (min)} & \textbf{Tipo} \\
\hline
01 & CH-001 & 8   & Nível 1 \\
02 & CH-002 & 12  & Nível 1 \\
03 & CH-003 & 180 & Nível 3 \\
04 & CH-004 & 9   & Nível 1 \\
05 & CH-005 & 15  & Nível 2 \\
06 & CH-006 & 11  & Nível 1 \\
07 & CH-007 & 240 & Nível 3 \\
08 & CH-008 & 10  & Nível 1 \\
09 & CH-009 & 14  & Nível 2 \\
10 & CH-010 & 13  & Nível 1 \\
\hline
\end{tabular}

\medskip
\textbf{Tarefas:}
\begin{enumerate}
  \item Calcule a média e a mediana do tempo total.
  \item Estime visualmente: se você desenhasse um histograma com bins
  $[0{-}20)$, $[20{-}100)$, $[100{-}200)$, $[200{-}300)$, quantas observações
  cairiam em cada bin?
  \item A distribuição tem assimetria positiva ou negativa? Como a média
  se compara à mediana?
  \item Qual medida central você reportaria para o SLA de atendimento?
\end{enumerate}

\textbf{Respostas esperadas:}
\begin{itemize}
  \item Soma: $8+12+180+9+15+11+240+10+14+13 = 512$ \quad \textbf{Média:} $51{,}2$ min
  \item Ordenados: $8,9,10,11,12,13,14,15,180,240$ \quad \textbf{Mediana:} $(12+13)/2 = 12{,}5$ min
  \item Bins: $[0{-}20)$: 8 chamados; $[100{-}200)$: 1; $[200{-}300)$: 1 — forte assimetria positiva
  \item Média (51,2) $\gg$ Mediana (12,5): média distorcida por Nível 3
  \item \textbf{Reportar:} Mediana = 12,5 min (honesto para Nível 1 e 2); investigar Nível 3 separadamente
\end{itemize}
\end{BoardBox}

% =============================================================================
\section{Parte III — Dados Ausentes (20h30–21h)}
% =============================================================================

\begin{SolvedBox}
\textbf{Por que dados ausentes importam?}

Dados ausentes não são apenas ``campos em branco''. Eles podem indicar:
\begin{itemize}
  \item Falha de sistema (campo que deveria ser preenchido mas não foi).
  \item Comportamento do respondente (pessoas omitem renda por vergonha).
  \item Processo seletivo (clientes que cancelam antes de completar pesquisa).
\end{itemize}

A forma como os dados estão ausentes muda \textbf{completamente} a estratégia correta.
\end{SolvedBox}

\begin{BoardBox}
\textbf{Os três tipos de missing (para o quadro)}

\textbf{1. MCAR — Missing Completely At Random (Ausente Completamente ao Acaso)}
\begin{itemize}
  \item A ausência não tem relação com nenhuma variável do dataset.
  \item Ex.: perda aleatória de conexão durante o preenchimento online.
  \item Estratégia: pode remover a linha sem introduzir viés significativo.
\end{itemize}

\textbf{2. MAR — Missing At Random (Ausente ao Acaso — mas relacionado a outra variável)}
\begin{itemize}
  \item A ausência está relacionada a outras variáveis observáveis, mas não à
  variável em si.
  \item Ex.: homens tendem a não responder perguntas sobre saúde mental em
  pesquisas — não por terem mais ou menos problemas, mas por gênero.
  \item Estratégia: imputar usando a variável relacionada (ex.: média do grupo de gênero).
\end{itemize}

\textbf{3. MNAR — Missing Not At Random (Ausente por razão relacionada ao próprio valor)}
\begin{itemize}
  \item A ausência está relacionada ao valor que deveria estar lá.
  \item Ex.: clientes insatisfeitos pulam a pergunta ``como foi sua experiência?'';
  funcionários não informam salário quando é muito alto ou muito baixo.
  \item Estratégia: \textbf{não impute sem entender o mecanismo}. Crie uma variável
  flag (\texttt{nota\_ausente = 1}) e analise o padrão de ausência separadamente.
  \item Este é o tipo mais perigoso — imputar leva a viés sistemático nas conclusões.
\end{itemize}
\end{BoardBox}

\begin{FormulaBox}{Estratégias de tratamento de dados ausentes}
\begin{tabular}{|l|p{5cm}|p{5cm}|}
\hline
\textbf{Estratégia} & \textbf{Quando usar} & \textbf{Risco} \\
\hline
Remover a linha & MCAR; $<\!5\%$ de ausências & Perde dados; viés se MAR/MNAR \\
Imputar pela média & Numérica simétrica, MCAR & Reduz variância; mascara distribuição \\
Imputar pela mediana & Numérica assimétrica ou com outliers & Idem, mas mais robusto \\
Imputar pela moda & Categórica & Pode criar dominância artificial \\
Criar flag + imputar & MNAR suspeito & Preserva o sinal da ausência \\
Imputação por grupo & MAR & Precisa da variável relacionada \\
\hline
\end{tabular}
\end{FormulaBox}

\begin{NoteBox}
\textbf{Documentação é obrigatória:} cada decisão sobre missing deve ser
registrada no relatório. Le leitor precisa saber: quantos valores foram
tratados, qual estratégia foi usada e por quê. Sem documentação, o dado
tratado é suspeito.
\end{NoteBox}

\begin{SolvedBox}
\textbf{Questão para o Quadro — Classificando Tipos de Missing}

Um hospital registrou dados de pacientes em pesquisa pós-operatória.
Classifique cada caso como \textbf{MCAR, MAR ou MNAR} e proponha a estratégia:

\begin{enumerate}[label=\alph*)]
  \item 5 de 300 formulários foram perdidos por falha no servidor.
  \item Pacientes idosos ($>$70 anos) tendem a não responder por
  dificuldade com o formulário digital.
  \item Pacientes com complicações pós-operatórias evitam avaliar a
  experiência cirúrgica.
\end{enumerate}

\textbf{Resolução:}
\begin{enumerate}[label=\alph*)]
  \item \textbf{MCAR} — perda aleatória (falha de sistema), sem relação com
  variáveis do paciente. Com 1,7\%, pode remover as linhas.
  \item \textbf{MAR} — ausência relacionada a outra variável observável
  (idade), não ao valor da resposta. Imputar pela mediana do grupo $>$70.
  \item \textbf{MNAR} — ausência ligada ao que seria respondido (avaliação
  negativa). \textbf{Não impute!} Criar flag \texttt{avaliacao\_ausente = 1}
  e investigar separadamente.
\end{enumerate}
\end{SolvedBox}

% =============================================================================
\section{Parte IV — Prática em Dupla 2 (21h–21h30)}
% =============================================================================

\begin{BoardBox}
\textbf{Tabela 2 — Pesquisa de Satisfação de Funcionários (com valores ausentes)}

\medskip
\begin{tabular}{|c|l|r|l|r|}
\hline
\textbf{ID} & \textbf{Nome} & \textbf{Nota (1–10)} & \textbf{Dept} & \textbf{Salário (R\$)} \\
\hline
01 & Ana     & 8    & TI      & 6.200 \\
02 & Bruno   & ---  & RH      & 4.800 \\
03 & Carla   & 7    & TI      & 6.100 \\
04 & Diego   & 2    & Vendas  & 3.200 \\
05 & Elena   & 9    & TI      & 7.800 \\
06 & Felipe  & ---  & Vendas  & ---   \\
07 & Giovana & 6    & RH      & 4.500 \\
08 & Hugo    & 8    & TI      & 6.400 \\
09 & Isabela & ---  & Vendas  & 3.100 \\
10 & João    & 5    & RH      & 4.600 \\
\hline
\end{tabular}

\medskip
\textbf{Tarefas:}
\begin{enumerate}
  \item Quantos valores ausentes há em cada coluna? Qual a porcentagem?
  \item Para Nota: qual tipo de missing você suspeita (MCAR, MAR ou MNAR)?
  Justifique olhando para o departamento e nota de Diego.
  \item Para Salário: o que a ausência do salário de Felipe sugere?
  \item Proponha uma estratégia de imputação para Nota e documente-a.
  \item Calcule a média de Nota antes e depois da imputação — mudou?
\end{enumerate}

\textbf{Respostas esperadas:}
\begin{itemize}
  \item Nota: 3 ausências (Bruno, Felipe, Isabela) — 30\%
  \item Salário: 1 ausência (Felipe) — 10\%
  \item Nota de Vendas: Diego = 2. Bruno e Isabela também são de Vendas (RH) e
  Vendas respectivamente. \textbf{Suspeita de MNAR:} insatisfeitos omitiram a resposta.
  Imputar pela média geral (7{,}0) criaria viés — subestimaria a insatisfação do Vendas.
  Estratégia correta: \textbf{criar flag} (\texttt{nota\_ausente = 1}) e imputar
  pela mediana do departamento Vendas (= 2).
  \item Salário de Felipe: ausência isolada, possivelmente MCAR (erro de sistema).
  Imputar pela mediana de Vendas (média entre 3.100 e 3.200 = 3.150).
  \item Média de Nota (sem ausentes): $(8+7+2+9+6+8+5)/7 = 6{,}43$
  \item Com imputação ingênua (média geral = 6,43 para os 3 ausentes): média geral
  permanece 6,43 — mas esconde a realidade de Vendas.
  \item Com imputação por grupo (Vendas = 2): média geral $\approx 5{,}7$ — mais honesta.
\end{itemize}
\end{BoardBox}

% =============================================================================
\section{Parte V — Preparação de Dados para Visualização (21h30–21h45)}
% =============================================================================

\begin{BoardBox}
\textbf{Checklist de preparação de dados para Tableau/Power BI (para o quadro)}

Antes de importar qualquer dataset em uma ferramenta de visualização:

\begin{tabular}{|p{4cm}|p{9cm}|}
\hline
\textbf{Ponto de verificação} & \textbf{O que fazer} \\
\hline
Tipos de dados corretos & Números como número (não texto). Datas como data (não texto). Categorias sem espaços extras.\\
Nomes de colunas limpos & Sem acento, sem espaço, sem caractere especial: \texttt{tempo\_min} não \texttt{Tempo (min)} \\
Datas no formato padrão & \texttt{YYYY-MM-DD} é o mais seguro internacionalmente. Evitar \texttt{DD/MM/AAAA} em inglês \\
Nulos tratados ou explícitos & Decidir: remover, imputar ou manter (com flag). Nunca deixar `?` ou `` `0` '' como proxy de nulo. \\
Unidades explícitas & Separar valor e unidade em colunas distintas: \texttt{valor\_reais}, \texttt{valor\_dolar} \\
Categorias padronizadas & ``SP'', ``São Paulo'', ``S.Paulo'' são três categorias diferentes para a ferramenta \\
\hline
\end{tabular}
\end{BoardBox}

\begin{NoteBox}
\textbf{Regra de ouro para Python/Pandas:} nunca deixe ``limpeza'' para o
último momento. Cada vez que você modifica um dataset sem documentação,
cria um arquivo que ninguém mais consegue reproduzir.\\
Pipeline recomendado: carrega $\to$ inspeciona $\to$ limpa $\to$ salva ($\to$ documenta).
\end{NoteBox}

\begin{SolvedBox}
\textbf{Questão para o Quadro — Diagnóstico de Qualidade de Dados}

Você recebeu um CSV e encontrou os seguintes problemas. Para cada um,
indique a correção:

\begin{center}
\begin{tabular}{|p{5.5cm}|p{7cm}|}
\hline
\textbf{Problema} & \textbf{Correção} \\
\hline
Coluna ``receita'' com texto ``R\$~3.500'' & Remover ``R\$'', converter para numérico \\
\hline
Datas: ``15/03/2026'' e ``2026-03-15'' & Padronizar: \texttt{YYYY-MM-DD} \\
\hline
Cidade: ``Recife'', ``recife'', ``RECIFE'' & \texttt{.str.title()} $\to$ ``Recife'' \\
\hline
Campo ``idade'' com valor $-5$ & Erro de digitação? Tratar como missing + documentar \\
\hline
\end{tabular}
\end{center}

\textbf{Regra:} cada problema encontrado deve ser documentado \emph{antes}
de ser corrigido. O relatório deve registrar: \textit{o quê} foi
encontrado, \textit{quantos} registros afetados, e \textit{como} foi tratado.
\end{SolvedBox}

% =============================================================================
\section{Parte VI — Atividade Individual Final (21h45–22h)}
% =============================================================================

\begin{SolvedBox}
\textbf{Instruções:}
\begin{enumerate}
  \item Calcule a média e mediana dos dados e compare-as.
  \item Identifique o tipo de assimetria (olhe para qual direção a cauda aponta).
  \item Identifique os missing values. Classifique cada um como MCAR, MAR ou MNAR.
  \item Proponha uma estratégia de imputação e calcule o valor imputado.
  \item Redija 1 insight: \emph{o dado mostra X, isso significa Y, portanto Z.}
\end{enumerate}
\textbf{Tempo:} 15 minutos. Entrega: folha ou foto no AVA.
\end{SolvedBox}

\begin{BoardBox}
\textbf{Tabela 3 — Tempo de Resposta de Proposta Comercial (em horas)}

\medskip
\begin{tabular}{|c|l|r|l|}
\hline
\textbf{ID} & \textbf{Consultor} & \textbf{Tempo (h)} & \textbf{Porte do cliente} \\
\hline
01 & Rafael   & 2   & PME \\
02 & Simone   & 4   & PME \\
03 & Tiago    & --- & Enterprise \\
04 & Ursula   & 3   & PME \\
05 & Vitor    & 2   & PME \\
06 & Wanda    & 96  & Enterprise \\
07 & Xande    & 3   & PME \\
08 & Yasmin   & 5   & PME \\
09 & Zé       & --- & Enterprise \\
10 & Amábile  & 4   & PME \\
\hline
\end{tabular}
\end{BoardBox}

\begin{SolvedBox}
\textbf{Gabarito (professor):}
\begin{itemize}
  \item Valores presentes: $2+4+3+2+96+3+5+4 = 119$ h; $n=8$; \textbf{Média:} $14{,}9$ h
  \item Ordenados: $2,2,3,3,4,4,5,96$ \quad \textbf{Mediana:} $(3+4)/2 = 3{,}5$ h
  \item Média (14,9) $\gg$ Mediana (3,5) $\Rightarrow$ assimetria positiva intensa (causada por Wanda/96h)
  \item \textbf{Missing de Tiago e Zé:} ambos Enterprise. Suspeita de \textbf{MNAR} — respostas
  a Enterprise podem demorar tanto que nunca foram registradas (ou o consultor evitou reportar).
  Estratégia: criar flag \texttt{atrasado\_enterprise = 1} + imputar pela mediana de Enterprise
  (que seria 96h, o único valor disponível).
  \item \textbf{Insight esperado:} ``O dado mostra que o tempo mediano de resposta para
  PME é 3,5h — dentro do SLA. Porém, clientes Enterprise têm ausência de dado (provável
  atraso não reportado) e o único caso registrado (96h) é 27× maior que a mediana PME.
  Portanto, recomendo investigar o processo de proposta para Enterprise antes de reportar
  uma taxa de cumprimento de SLA que ignora esses casos.''
\end{itemize}
\end{SolvedBox}

% =============================================================================
\section{Revisão Semanal — Questão Tipo Prova I}
% =============================================================================

\begin{SolvedBox}
\textbf{Questão de Revisão (Semana 05) — Distribuições, Histogramas e Dados Ausentes}

Uma fintech que oferece microcrédito para empreendedores do Nordeste está
analisando sua base de 15.000 clientes. As variáveis principais são:

\begin{center}
\begin{tabular}{|l|l|c|l|}
\hline
\textbf{Variável} & \textbf{Tipo} & \textbf{\% Missing} & \textbf{Distribuição} \\
\hline
\texttt{renda\_mensal}       & Numérica    & 8\%  & Assimétrica à direita \\
\texttt{tempo\_negocio}      & Numérica    & 15\% & Bimodal (picos em 6 e 36 meses) \\
\texttt{segmento}            & Categórica  & 22\% & Alimentação 40\%, Vestuário 25\%, Serviços 35\% \\
\hline
\end{tabular}
\end{center}

A analista construiu histogramas e observou:
\begin{itemize}
  \item \texttt{renda\_mensal}: maioria entre R\$1.000 e R\$3.000, cauda até R\$25.000.
  \item \texttt{tempo\_negocio}: dois picos — 6 meses (iniciantes) e 36 meses (estabelecidos).
\end{itemize}

O gerente argumenta que os 22\% de missing em \texttt{segmento} são problemáticos
porque essa variável é essencial para segmentar dashboards.

\textbf{Pergunta:} Qual medida é mais representativa para \texttt{renda\_mensal}?
O que significa a bimodalidade de \texttt{tempo\_negocio}? Qual a melhor
estratégia para os 22\% de missing em \texttt{segmento}?

\medskip
\textbf{Resolução:}

\begin{enumerate}
  \item \textbf{Renda:} distribuição assimétrica à direita $\Rightarrow$ média $>$ mediana.
  A \textbf{mediana} é mais representativa porque não é distorcida pela cauda de rendas altas.

  \item \textbf{Tempo de negócio:} a distribuição bimodal indica \textbf{dois subgrupos distintos}
  — empreendedores iniciantes ($\approx$ 6 meses) e estabelecidos ($\approx$ 36 meses).
  Analisar conjuntamente mascara ambos; o correto é segmentar.

  \item \textbf{Segmento:} criar uma categoria \textbf{``Não informado''} e incluí-la
  como filtro no Tableau, permitindo análises com e sem esses registros.
  Imputar com a moda (``Alimentação'') distorceria a distribuição real.
  Remover 22\% dos registros descartaria dados valiosos das demais variáveis.
\end{enumerate}
\end{SolvedBox}

% ---------------------------------------------------------------------------
\section{Materiais Preparados}
% ---------------------------------------------------------------------------
\begin{itemize}
  \item Slides com animação: histograma construído bin a bin
  \item Galeria de distribuições (normal, positiva, negativa, bimodal)
  \item Tabelas 1, 2 e 3 impressas
  \item Notebook base (\texttt{semana05\_distribuicoes.ipynb}) com dados carregados
\end{itemize}

% ---------------------------------------------------------------------------
\section{Próximo passo}
% ---------------------------------------------------------------------------

Na \textbf{Semana 06} chegam as apresentações do Trabalho 3 e entramos
na análise \textbf{bivariada}: o que acontece quando duas variáveis numéricas
se movem juntas? Correlação de Pearson, scatter plots e — o tema mais
mal interpretado em dados — a diferença entre correlação e causalidade.

Perguntem aos alunos: \textit{``Se eu te disser que países com mais
televisores por domicílio têm maior expectativa de vida, o que você
concluiría? Faz sentido distribuir televisores para aumentar a expectativa
de vida? O que essa correlação realmente está capturando?''}
